\begin{lemma}[Situation 7.1: advance a stored common color]
Fix a replacement presentation \(P\) and
\(1\leq i<\ell-1\).  Suppose
\(\alpha\in D_{n,q}\) is missing under its coloring at both exact vertices
\(y_i\) and \(y_\ell\).  Under
\(\noderef{HierarchyPathNumber}=1\) and the explicit palette-surplus
hypothesis, one has
\[
 i+1\in\noderef{HierarchyEligibleCommonMissingIndices}.
\]
The conclusion exposes the replacement presentation obtained by the
paper's two consecutive Kempe changes.
\end{lemma}
\begin{proof}
Put \(\alpha=\eta_h\) and write
\(\Gamma^q_{\eta_h}=\{\gamma_1,\gamma_2\}\).  These are two distinct exact
colors by the cardinality field of \noderef{GoodHierarchy}.  The path-number
one hypothesis says that the exact edge \(e_{i+1}\) joins
\(y_i\) to \(y_{i+1}\).  At most one of the two reserved colors is its
\(\mu_n\)-color, so rename the pair so that
\[
             \mu_n(e_{i+1})\ne\gamma_1.                      \tag{2}
\]

Elementarity of the exact penultimate prefix follows from
\noderef{HierarchyCounterexamplePenultimateElementary} applied to the
transported minimum context in
\noderef{HierarchyReplacementPresentation}.  Since \(\eta_h\) is missing
at \(y_i\), it is not missing earlier in that prefix.  The avoidance field
of \noderef{GoodHierarchy} says that \(\gamma_1\) is absent from every
exact displayed tail segment before the first \(\eta_h\)-missing position.
Together with (2) and
\noderef{HierarchyTailMissingColorExclusion}, this gives the paper's exact
edge assertion
\[
 \gamma_1\notin\mu_n\langle T(y_{i+1})-A\rangle,
\]
and says that \(\eta_h\) is absent there except possibly at the certified
zero-divider SE/RE first seed.  In that exception the actual individual
seed identity and its membership in \(D_n\) are retained.

By \noderef{ChangeLemma}, the
\((\eta_h,\gamma_1)\)-component through the first missing witness \(y_i\)
is disjoint from the component through \(y_\ell\).  Swap the latter and
call the result \(\nu_n\).  The changed component contains no exact tail
edge: ordinary augmentation edges with either affected color lie in the
component through their earlier missing witness, and the sole SE/RE seed
is likewise in that component by its stored-color certificate.  Thus all
displayed tail-edge values agree.  Component separation makes all vertex
missing sets agree through
\noderef{HierarchyFullTailPrefix}\((C,\ell-1)\).  All unaffected colors
have the same candidate set for
\noderef{HierarchyFirstMissingPosition}; \(\eta_h\) retains its witness
\(y_i\), and \(\gamma_1\) retains its earlier hierarchy-stage witness.
These prove all three pointwise transport premises, not an equality of
whole terminal missing sets.

Apply \noderef{StableChangeCasesOneTwoFullTail}.  Its Cases 1/2 location
alternative is supplied by the stage-missing property of \(\gamma_1\).
It gives a full corresponding ETT, maximum property, a rebuilt base and
good hierarchy, and \noderef{SameGoodHierarchy} data.  Rung agreement
rewrites the lower-rung and fewer-divider functions.
\noderef{SameGoodHierarchy} identifies divider counts, stages, and hence
\noderef{HierarchyCounterexampleMeasure}, transferring lexicographic
minimality.  Corresponding ladder equality and stability transfer the
\(2n+11\) estimate, while the corresponding certificate supplies
properness and exact one-uncoloredness.  Thus every field of
\noderef{HierarchyMinimumCounterexampleContext} is reconstructed and the
outputs form a replacement presentation \(P_1\).  Under \(P_1\), \(\gamma_1\) is missing
at \(y_\ell\) and at the last ordinary stage.

Choose \(\theta\in\overline{\nu_n}(y_{i+1})\) using the palette-surplus
hypothesis.  Elementarity of the unchanged penultimate prefix gives
\(\theta\ne\gamma_1\) and says that \(\theta\) has no earlier missing
witness.  The exact-edge exclusion above and (2) imply, via
\noderef{HierarchyTailMissingColorExclusion}, that \(\theta\) occurs on no
edge through \(y_{i+1}\), apart from the explicitly certified first-seed
exception.

Apply \noderef{ChangeLemma} to \(\theta,\gamma_1\).  Its first component
ends at \(y_{i+1}\) and is disjoint from the component through \(y_\ell\).
Swap the latter.  Exactly the same augmentation/seed dichotomy proves
tail-edge equality, component separation proves penultimate vertexwise
missing-set equality, and the earlier witnesses prove stored
first-position equality.  Therefore a second application of
\noderef{StableChangeCasesOneTwoFullTail} constructs a replacement
presentation \(P_2\).  The resulting coloring leaves \(\theta\) missing at
\(y_{i+1}\) and makes it missing at \(y_\ell\).  Since
\(1\leq i+1\leq\ell-1\), \(i+1,P_2,\theta\) satisfy the filter defining
\noderef{HierarchyEligibleCommonMissingIndices}.
\end{proof}
